% Generated draft: requires mathematical and source review.
\documentclass[11pt]{article}
\usepackage[margin=25mm]{geometry}
\usepackage{fontspec}
\setmainfont{Latin Modern Roman}
\usepackage{amsmath,amssymb,mathtools,unicode-math}
\setmathfont{Latin Modern Math}
\usepackage{enumitem,needspace,xurl,hyperref}
\setlength{\parindent}{0pt}
\setlength{\parskip}{5pt}
\setlength{\emergencystretch}{3em}
\newcommand{\N}{\mathbb{N}}
\newcommand{\Z}{\mathbb{Z}}
\newcommand{\Q}{\mathbb{Q}}
\newcommand{\R}{\mathbb{R}}
\newcommand{\C}{\mathbb{C}}
\newcommand{\F}{\mathbb{F}}
\newcommand{\A}{\mathbb{A}}
\newcommand{\PP}{\mathbb{P}}
\newcommand{\E}{\mathbb{E}}
\newcommand{\eps}{\varepsilon}
\newcommand{\dd}{\,\mathrm{d}}
\begin{document}
\section*{UnsolvedMath ID: 801057}
\textbf{Source problem code:} catalogue entry
\subsection{Definitions and mathematical statement}
An odd period is a positive integer of the form \(2k+1\), where \(k\in\Z_{\geq 0}\).

The supplied source asks the following existence question: does there exist a mildly dissipative map \(f_{c,b}\in\mathcal E_0\) that admits arbitrarily deep renormalizations with odd periods?

Here ``arbitrarily deep'' requires renormalizations at depths exceeding every prescribed finite bound. The source also motivates the question by stating that it is unknown whether different combinatorics of renormalization may occur; this is a supplied contextual assessment.

The supplied record does not define the family \(f_{c,b}\), its parameter domain, the class \(\mathcal E_0\), mild dissipativity, or renormalization and its depth. It also does not specify whether an odd period is measured for the original map or relative to the preceding renormalization, or whether the required renormalizations must belong to one nested sequence. These source-dependent notions must be recovered before this question has a fully specified standalone mathematical formulation.
\subsection{Short English statement}
The question asks whether some mildly dissipative map in the source's class \(\mathcal E_0\) has odd-period renormalizations at arbitrarily large depths. Essential definitions and the intended period convention are missing from the supplied record.
\subsection{Sources}
\begin{itemize}
\item Supplied local UnsolvedMath registry record, permanent ID \(801057\), exact source problem code \texttt{catalogue entry}, titled ``Henon Maps: Open Problems --- Question 9.'' The record attributes the problem to Berger and Bedford, with proposed year 2023, and labels it open. This status is supplied metadata, not an independently verified mathematical or literature assessment.
\item Berger and Bedford, \emph{Henon Maps: A List of Open Problems} (2023), Question 9, \url{https://arxiv.org/abs/2312.03907}. The supplied provenance records access on 2026-10-16 and extraction by \texttt{arxiv\_latex}; these are source-provided provenance claims, not a claim of independent consultation.
\end{itemize}
\end{document}
